\section{Introduction}
\label{sec:intro}

Modern LLM agents re-read long, mostly-static instructions on every turn: a system policy, a tool
specification, retrieved documents. Key/value (KV) caching makes this affordable by reusing the prefill
computation across turns---but only across an \emph{exact} shared prefix. The moment a single token
changes inside the reused region---a timestamp, a user id, an order's status---the keys and values of
\emph{every} subsequent token are invalidated, because each attended to the token that changed. The
de-facto workaround is to \emph{hoist} all mutable content to the end of the prompt so the static prefix
stays cache-aligned. This pushes an inference-layer constraint into the application layer: fields
referenced in several places, nested sub-agent prompts, and dynamically assembled contexts cannot all be
cleanly hoisted, and the application must enumerate every mutable field in advance.

This paper starts from a concrete puzzle. The region of the cache \emph{before} a field is, by
construction, independent of the field's value---we measure a key/value deviation of exactly $0.0$ when
the field changes. One might therefore hope to \emph{surgically} refresh only the field's own keys and
values, leave the rest of the cache stale, and pay almost nothing. We find this fails completely: the
model's decision reverts to the \emph{old} field value, as if the edit never happened
(\Cref{fig:teaser}b).

\paragraph{The discovery.} The reason, which we establish causally, is that transformers do not defer
their reasoning to decode time. At \emph{prefill} the model already computes the
\emph{field-conditioned conclusion} and writes it onto downstream tokens---disproportionately onto
aggregator/delimiter tokens that later positions attend through. The decision then reads these
\emph{notes}, not the field itself: across models the field's own KV causally drives less than $1\%$ of
the decision, while the downstream notes drive essentially all of it. The KV cache is best understood not
as a frozen byproduct of prefill but as a \emph{notebook of memoized conclusions} (\Cref{fig:teaser}a).

\paragraph{Two capabilities from one mechanism.} This reframing is generative. (1) If the conclusion is
already written downstream, then \emph{editing} a field means amending the notes, not recomputing them:
a one-line salient \emph{erratum} suffices, and we map exactly when the cheap in-place edit alone is
enough (under reasoning) versus when it is not. (2) If the notes are localized and position-portable,
then a reusable skill can be \emph{composed} into a new context by repositioning and splicing its cached
notes---no recompute. A \emph{keystone} experiment, editing a field \emph{inside} a transplanted skill,
shows the two operations act on a single substrate (\Cref{fig:overview}).

\paragraph{Contributions.}
\begin{itemize}[leftmargin=1.2em,itemsep=1pt,topsep=2pt]
\item \textbf{A mechanism} (\Cref{sec:mechanism}): \emph{attention-mediated memoized inference}, established
by four independent causal probes (locality patching, suffix concentration, linear probing, circuit
knockout), connecting to delimiter-token aggregation observed in interpretability.
\item \textbf{An editing capability} (\Cref{sec:editable}): naive KV editing fails; the erratum /
\texttt{field+erratum} fix matches the hoist-to-end oracle without prompt surgery; a reasoning/scale
analysis of when the ${\sim}1\%$-compute in-place edit suffices.
\item \textbf{A composing capability} (\Cref{sec:composable}): position-portable transplant of precompiled
skills, $O(L)$ vs.\ $O(L^2)$ TTFT ($13.9\times$ at $32$k), with a seam-repair knob---honestly credited to
prior caching work, with our addition being the mechanism that explains it and a correctness lens.
\item \textbf{The unification} (\Cref{sec:keystone}): the keystone and a unified edit+compose agent over
twelve models.
\item \textbf{Reach and scope} (\Cref{sec:reach}): validation across scale, MoE, quantization, and
multimodal image caches; small adapters for MLA and interleaved M-RoPE; a fix for sliding-window; and an
analysis of where the substrate holds up to the 2026 sparse/compressed-attention frontier.
\item \textbf{Systems payoff} (\Cref{sec:systems}): a real agentic environment and a closed vLLM
integration ($16\times$ throughput).
\item \textbf{An application: user memory} (\Cref{sec:memory}): the large, mutable user-memory
document is both composed (precompiled, repositioned, spliced) and edited (in-place / erratum)
as one substrate---decision-faithful to full recompute at $2.3$--$4.25\times$ lower
time-to-first-token, validated to $70$B and on real long-conversation memory (LoCoMo, transplant
$\equiv$ full recompute in QA accuracy)---with a pre-registered, statistically-controlled evaluation.
\end{itemize}

\begin{figure}[t]
\centering
\includegraphics[width=\linewidth]{fig1_teaser}
\caption{\textbf{Models take notes at prefill.} (a) At prefill the model memoizes the field-conditioned
conclusion onto downstream aggregator tokens (orange); at decode the decision reads those notes (blue).
(b) Consequently, surgically editing the field's own KV is ignored---the fix is an \emph{erratum} that
amends the notes. (c) Because the notes are position-portable, a precompiled skill can be pasted into a
new context in $O(L)$ rather than $O(L^2)$ time.}
\label{fig:teaser}
\end{figure}

\begin{figure}[t]
\centering
\includegraphics[width=0.92\linewidth]{fig0_overview}
\caption{\textbf{One mechanism, two capabilities, one substrate.} The discovery that models memoize
field-conditioned conclusions at prefill makes the KV cache both \emph{editable} (amend the notes with a
salient erratum) and \emph{composable} (reposition and splice precompiled notes); a keystone experiment
shows they act on a single substrate, which holds up to the 2026 attention frontier.}
\label{fig:overview}
\end{figure}
